\chapter{Cohomology of Sheaves}

\section{The Global Section Functor}

\begin{definition}
Let \(X\) be a topological space and let \(\FF\) be a sheaf of abelian groups.
The group of \emph{global sections} is
\[
    \Gamma(X,\FF)=\FF(X).
\]
This defines a functor \(\Gamma(X,-)\) from sheaves of abelian groups on \(X\)
to abelian groups.
\end{definition}

\begin{proposition}
The global section functor is left exact: if
\[
    0\longrightarrow \FF'\longrightarrow \FF\longrightarrow \FF''
\]
is exact, then
\[
    0\longrightarrow \Gamma(X,\FF')\longrightarrow \Gamma(X,\FF)
    \longrightarrow \Gamma(X,\FF'')
\]
is exact.
\end{proposition}

\begin{proof}
Injectivity of \(\Gamma(X,\FF')\to\Gamma(X,\FF)\) follows from injectivity on
stalks and the fact that sections are determined by germs.  If
\(s\in\Gamma(X,\FF)\) maps to zero in \(\Gamma(X,\FF'')\), then each germ
\(s_x\) lies in the image of \(\FF'_x\to\FF_x\).  Hence locally \(s\) lifts to a
section of \(\FF'\).  These local lifts agree on overlaps because
\(\FF'\to\FF\) is injective, so they glue to a global section of \(\FF'\)
mapping to \(s\).  This proves exactness at \(\Gamma(X,\FF)\).
\end{proof}

\section{Flasque Sheaves and Sheaf Cohomology}

\begin{definition}
A sheaf \(\FF\) on \(X\) is \emph{flasque} if for every inclusion of open sets
\(V\subset U\), the restriction map \(\FF(U)\to\FF(V)\) is surjective.
\end{definition}

\begin{lemma}
Let
\[
    0\longrightarrow \FF'\longrightarrow \FF\longrightarrow \FF''
    \longrightarrow 0
\]
be an exact sequence of sheaves of abelian groups.  If \(\FF'\) is flasque, then
\[
    0\longrightarrow \Gamma(X,\FF')\longrightarrow \Gamma(X,\FF)
    \longrightarrow \Gamma(X,\FF'')\longrightarrow 0
\]
is exact.
\end{lemma}

\begin{proof}
Only surjectivity on the right is not already contained in left exactness.  Let
\(s''\in\Gamma(X,\FF'')\).  Choose an open cover \(\{U_i\}\) and lifts
\(s_i\in\FF(U_i)\) of \(s''|_{U_i}\).  We refine the cover to a well ordered
family and construct compatible lifts by transfinite induction.  Suppose a lift
has been constructed on an open set \(V\), and let \(U_i\) be the next open set.
On \(V\cap U_i\), the difference between the existing lift and \(s_i\) maps to
zero in \(\FF''\), hence comes from a section \(a\in\FF'(V\cap U_i)\).  Since
\(\FF'\) is flasque, extend \(a\) to a section \(b\in\FF'(U_i)\).  Replacing
\(s_i\) by \(s_i+b\), the two lifts agree on \(V\cap U_i\), and hence glue to a
lift over \(V\cup U_i\).  At limit stages the previously constructed compatible
lifts glue by the sheaf axiom.  The final result is a global lift of \(s''\).
\end{proof}

\begin{definition}
For a sheaf \(\FF\) of abelian groups on \(X\), define
\[
    C^0(\FF)(U)=\prod_{x\in U}\FF_x.
\]
The natural map \(\FF\to C^0(\FF)\) sends a section to its family of germs.  Let
\(Q^0(\FF)=\coker(\FF\to C^0(\FF))\), and define inductively
\[
    C^n(\FF)=C^0(Q^{n-1}(\FF)),\qquad
    Q^n(\FF)=\coker(Q^{n-1}(\FF)\to C^n(\FF)).
\]
The resulting complex
\[
    0\longrightarrow \FF\longrightarrow C^0(\FF)\longrightarrow C^1(\FF)
    \longrightarrow C^2(\FF)\longrightarrow\cdots
\]
is the \emph{Godement resolution} of \(\FF\).
\end{definition}

\begin{proposition}
Each \(C^n(\FF)\) is flasque, and the Godement resolution is exact.
\end{proposition}

\begin{proof}
For \(C^0(\FF)\), restriction from \(U\) to \(V\) is the projection
\(\prod_{x\in U}\FF_x\to\prod_{x\in V}\FF_x\), which is surjective.  Hence
\(C^0(\FF)\) is flasque, and the same argument applies to \(C^n(\FF)\) because
each is \(C^0\) of another sheaf.

The map \(\FF\to C^0(\FF)\) is injective because a section whose germs are all
zero is zero.  By definition \(Q^0(\FF)\) is its cokernel, so
\[
    0\to \FF\to C^0(\FF)\to Q^0(\FF)\to 0
\]
is exact.  Repeating the same construction for \(Q^0(\FF)\), then for
\(Q^1(\FF)\), and so on, gives exactness at every term of the displayed
resolution.
\end{proof}

\begin{definition}
The \emph{sheaf cohomology groups} of \(\FF\) on \(X\) are
\[
    H^i(X,\FF)=H^i\bigl(\Gamma(X,C^\bullet(\FF))\bigr),
\]
where \(C^\bullet(\FF)\) is the Godement resolution without the initial
\(\FF\)-term.  Thus \(H^0(X,\FF)=\Gamma(X,\FF)\).
\end{definition}

\begin{theorem}[Long exact sequence]
Every short exact sequence of sheaves of abelian groups
\[
    0\longrightarrow \FF'\longrightarrow \FF\longrightarrow \FF''
    \longrightarrow 0
\]
induces a natural long exact sequence
\[
\begin{aligned}
0&\to H^0(X,\FF')\to H^0(X,\FF)\to H^0(X,\FF'')\\
&\to H^1(X,\FF')\to H^1(X,\FF)\to H^1(X,\FF'')\to\cdots .
\end{aligned}
\]
\end{theorem}

\begin{proof}
The functor \(C^0\) is exact: a short exact sequence of sheaves gives short
exact sequences on stalks, and products of short exact sequences of abelian
groups remain exact.  Applying the construction of \(C^\bullet\) inductively
therefore gives a short exact sequence of Godement resolutions
\[
    0\to C^\bullet(\FF')\to C^\bullet(\FF)\to C^\bullet(\FF'')\to0.
\]
The kernels in each degree are flasque, so the previous lemma shows that
applying \(\Gamma(X,-)\) preserves exactness degree by degree.  Thus we obtain
a short exact sequence of complexes of abelian groups.  The usual snake-lemma
construction for a short exact sequence of complexes gives the connecting
homomorphisms and the long exact cohomology sequence.
\end{proof}

\section{\texorpdfstring{\v{C}ech}{Cech} Cohomology}

\begin{definition}
Let \(\UU=\{U_i\}_{i\in I}\) be an open cover of \(X\), with \(I\) totally
ordered.  For a sheaf of abelian groups \(\FF\), define the \emph{\v{C}ech
cochains}
\[
    \cech^p(\UU,\FF)
    =
    \prod_{i_0<\cdots<i_p}
    \FF(U_{i_0}\cap\cdots\cap U_{i_p}).
\]
The differential \(d:\cech^p(\UU,\FF)\to\cech^{p+1}(\UU,\FF)\) is
\[
    (dc)_{i_0\ldots i_{p+1}}
    =
    \sum_{a=0}^{p+1}(-1)^a
    c_{i_0\ldots\widehat{i_a}\ldots i_{p+1}}
    \big|_{U_{i_0}\cap\cdots\cap U_{i_{p+1}}}.
\]
The cohomology of this complex is denoted \(\cH^p(\UU,\FF)\).
\end{definition}

\begin{proposition}
The \v{C}ech differential satisfies \(d^2=0\).
\end{proposition}

\begin{proof}
For a \(p\)-cochain \(c\), the component \((d^2c)_{i_0\ldots i_{p+2}}\) is a
sum over pairs \(a<b\).  The term obtained by first deleting \(i_b\) and then
deleting \(i_a\) has sign \((-1)^b(-1)^a\).  The same restricted section is
obtained by first deleting \(i_a\) and then deleting \(i_b\), but in the second
step the index \(b\) has shifted to \(b-1\), so the sign is
\((-1)^a(-1)^{b-1}\).  These two signs are opposite.  All terms cancel in
pairs, hence \(d^2=0\).
\end{proof}

\begin{proposition}
\(\cH^0(\UU,\FF)\simeq\Gamma(X,\FF)\).
\end{proposition}

\begin{proof}
A \(0\)-cocycle is a family \(s_i\in\FF(U_i)\) such that
\[
    s_i|_{U_i\cap U_j}=s_j|_{U_i\cap U_j}
\]
for all \(i,j\).  By the sheaf axiom, such a family glues uniquely to a global
section of \(\FF\).  This identifies \(0\)-cocycles with \(\Gamma(X,\FF)\), and
there are no negative cochains, so \(\cH^0=\ker d^0\).
\end{proof}

\section{Affine Acyclicity}

\begin{lemma}[Affine \v{C}ech exactness]
Let \(A\) be a ring, \(M\) an \(A\)-module, and let \(f_1,\ldots,f_r\in A\)
generate the unit ideal.  The augmented complex
\[
0\to M\to \prod_i M_{f_i}\to \prod_{i<j}M_{f_if_j}\to\cdots
\to M_{f_1\cdots f_r}\to0
\]
with the \v{C}ech differential is exact.
\end{lemma}

\begin{proof}
Choose \(a_1,\ldots,a_r\in A\) with \(\sum_i a_i f_i=1\).  For a
\(p\)-cochain \(c=(c_{i_0\ldots i_p})\), define a homotopy
\[
    (hc)_{i_0\ldots i_{p-1}}
    =
    \sum_j a_j f_j\,
    c_{j i_0\ldots i_{p-1}}
\]
after restricting all terms to the common localization; if the indices are not
ordered, reorder them and insert the sign of the permutation.  A direct
calculation gives
\[
    dh+hd=\id
\]
in positive degrees, because the terms in which \(d\) deletes one of the
original indices cancel with the corresponding terms from \(hd\), and the terms
in which \(d\) deletes the inserted index \(j\) sum to
\(\sum_j a_jf_j c=c\).  In degree \(0\), the same identity says that every
cocycle in \(\prod_i M_{f_i}\) comes from \(M\).  Thus the augmented complex is
exact.
\end{proof}

\begin{theorem}[Comparison for acyclic covers]
Let \(\UU=\{U_i\}\) be an open cover of \(X\).  Suppose that for every finite
intersection \(U_{i_0}\cap\cdots\cap U_{i_p}\) and every \(q>0\),
\[
    H^q(U_{i_0}\cap\cdots\cap U_{i_p},\FF)=0.
\]
Then the natural map
\[
    \cH^p(\UU,\FF)\longrightarrow H^p(X,\FF)
\]
is an isomorphism for all \(p\ge0\).
\end{theorem}

\begin{proof}
Resolve \(\FF\) by its Godement resolution \(C^\bullet(\FF)\).  Form the double
complex
\[
    K^{p,q}=\cech^p(\UU,C^q(\FF)).
\]
Taking cohomology first in the \(q\)-direction gives the \v{C}ech complex
\(\cech^\bullet(\UU,\FF)\), because the augmented Godement resolution is exact
on each finite intersection by the assumed vanishing of higher cohomology
there.  Taking cohomology first in the \(p\)-direction gives
\(\Gamma(X,C^\bullet(\FF))\): the sheaves \(C^q(\FF)\) are flasque, and for a
flasque sheaf every \v{C}ech cocycle is a coboundary in positive degree by the
same extension argument used in the proof that flasque kernels make global
sections exact.  The two spectral sequences of a first-quadrant double complex
therefore have the same total cohomology, and the two descriptions of that
total cohomology give the asserted isomorphism.
\end{proof}

\begin{theorem}[Affine acyclicity]
Let \(X=\Spec A\) and let \(\FF=\widetilde M\) be a quasi-coherent sheaf.  Then
\[
    H^i(X,\FF)=0\qquad\text{for all }i>0.
\]
\end{theorem}

\begin{proof}
It is enough to compute cohomology with respect to a finite principal affine
cover, because every finite intersection of principal affine opens is again a
principal affine open and the comparison theorem applies inductively.  For a
cover \(X=\bigcup_i D(f_i)\) with \((f_1,\ldots,f_r)=A\), the \v{C}ech complex
of \(\widetilde M\) is exactly the localization complex in the affine
\v{C}ech exactness lemma:
\[
0\to M\to \prod_i M_{f_i}\to \prod_{i<j}M_{f_if_j}\to\cdots .
\]
The augmented complex is exact, so the positive \v{C}ech cohomology groups
vanish.  The comparison theorem identifies these groups with sheaf cohomology.
\end{proof}

\begin{corollary}
If \(X\) is separated and \(\UU=\{U_i\}\) is an affine open cover, then for every
quasi-coherent sheaf \(\FF\) the natural map
\[
    \cH^p(\UU,\FF)\longrightarrow H^p(X,\FF)
\]
is an isomorphism for all \(p\ge0\).
\end{corollary}

\begin{proof}
For a separated scheme, every finite intersection of affine opens is affine.
The restriction of a quasi-coherent sheaf to such an intersection is
quasi-coherent.  Affine acyclicity gives the vanishing hypothesis in the
comparison theorem.
\end{proof}

\section{Serre's Criterion and Finiteness}

\begin{theorem}[Serre's affine criterion]
Let \(X\) be a quasi-compact scheme with a basis of affine open subsets.  Then
the following are equivalent:
\begin{enumerate}[label=(\roman*)]
    \item \(X\) is affine;
    \item \(H^i(X,\FF)=0\) for every quasi-coherent sheaf \(\FF\) and every
    \(i>0\);
    \item \(H^1(X,\II)=0\) for every quasi-coherent ideal sheaf
    \(\II\subset\OO_X\).
\end{enumerate}
\end{theorem}

\begin{proof}
The implication (i) \(\Rightarrow\) (ii) is affine acyclicity, and
(ii) \(\Rightarrow\) (iii) is immediate.

Assume (iii).  Let \(x\in X\).  Choose an affine open neighbourhood \(U\) of
\(x\), and then a principal open \(D(f)\subset U\) containing \(x\).  Let
\(Z=X\setminus D(f)\), with the reduced closed subset structure locally on
affine opens, and let \(\II\) be the quasi-coherent ideal sheaf of functions
vanishing on \(Z\).  On \(U\), the function \(f\) is a section of \(\II|_U\).
The obstruction to extending \(f\) from \(U\) to a global section of some power
\(\II^n\) lies in \(H^1(X,\II^n)\), which is zero by (iii).  Thus, after
replacing \(f\) by a power, there is a global section \(g\in\Gamma(X,\OO_X)\)
whose non-vanishing locus \(X_g\) satisfies
\[
    x\in X_g\subset D(f)\subset U.
\]
Such opens \(X_g\) form a basis of \(X\).

Let \(A=\Gamma(X,\OO_X)\).  The global functions define a morphism
\[
    \varphi:X\longrightarrow \Spec A.
\]
For every basis open \(X_g\), the restriction map identifies
\[
    \Gamma(X_g,\OO_X)\simeq A_g;
\]
this follows by applying the vanishing of \(H^1\) to the exact sequences
\[
    0\to\OO_X \xto{g^n} \OO_X\to \OO_X/g^n\OO_X\to0
\]
and passing to the localization colimit.  Therefore
\(\varphi|_{X_g}:X_g\to D(g)\) is an isomorphism of affine schemes.  Since the
opens \(X_g\) cover \(X\) and the corresponding \(D(g)\) cover \(\Spec A\),
\(\varphi\) is an isomorphism.  Hence \(X\) is affine.
\end{proof}

\begin{theorem}[Serre finiteness for projective schemes]
Let \(k\) be a field, let \(X\) be a projective \(k\)-scheme, and let \(\FF\) be
a coherent sheaf on \(X\).  Then each \(H^i(X,\FF)\) is a finite-dimensional
\(k\)-vector space and \(H^i(X,\FF)=0\) for \(i>\dim X\).
\end{theorem}

\begin{proof}
Embed \(X\) as a closed subscheme of \(\PP^n_k\).  The standard affine cover of
\(\PP^n_k\) restricts to a finite affine open cover \(\UU\) of \(X\).  Since
\(X\) is separated, finite intersections in this cover are affine.  By the
comparison theorem, \(H^i(X,\FF)\) is computed by the finite \v{C}ech complex
\[
    \prod_{j_0<\cdots<j_p}
    \Gamma(U_{j_0}\cap\cdots\cap U_{j_p},\FF).
\]
On each affine intersection, \(\FF\) corresponds to a finitely generated algebra
module over a finitely generated \(k\)-algebra.  In a fixed cohomological degree
the \v{C}ech cohomology is a subquotient of a finite product of such modules.
For projective schemes, the homogeneous grading bounds the degrees contributing
to a fixed coherent sheaf on this finite standard cover; equivalently, after a
finite presentation by sums of twists \(\OO_{\PP^n}(m)\), the cohomology of
the twists is the standard finite-dimensional cohomology of projective space.
Thus every \(H^i(X,\FF)\) is finite-dimensional.  Since a finite affine cover of
a noetherian scheme of dimension \(d\) has no non-empty intersections of more
than \(d+1\) members after refinement by a dimension filtration, the Cech
complex may be chosen to have length \(d\), giving vanishing for \(i>d\).
\end{proof}

\begin{definition}
For a coherent sheaf \(\FF\) on a projective \(k\)-scheme \(X\), define
\[
    \chi(X,\FF)=\sum_i (-1)^i\dim_k H^i(X,\FF).
\]
The sum is finite by Serre finiteness.
\end{definition}

\begin{proposition}[Additivity of Euler characteristic]
Let \(X\) be a projective \(k\)-scheme and let
\[
    0\longrightarrow \FF'\longrightarrow \FF\longrightarrow \FF''
    \longrightarrow0
\]
be an exact sequence of coherent sheaves.  Then
\[
    \chi(X,\FF)=\chi(X,\FF')+\chi(X,\FF'').
\]
\end{proposition}

\begin{proof}
The short exact sequence gives a long exact sequence of finite-dimensional
\(k\)-vector spaces.  Break the long exact sequence into short exact sequences
by taking kernels and images at every term.  The alternating sum of dimensions
in an exact finite complex is zero, because each image is counted once with a
plus sign and once with a minus sign.  Applying this to the long exact
cohomology sequence gives the displayed equality.
\end{proof}

\section{Cohomology of Projective Space}

\begin{theorem}
Let \(A\) be a ring and let \(X=\PP^n_A\).  Then
\[
    H^0(X,\OO_X(m))\simeq A[x_0,\ldots,x_n]_m
\]
for \(m\ge0\), and \(H^0(X,\OO_X(m))=0\) for \(m<0\).  Moreover
\[
    H^q(X,\OO_X(m))=0
    \quad\text{for }0<q<n
\]
for all \(m\), and \(H^n(X,\OO_X(m))\) is generated by Cech classes represented
by Laurent monomials
\[
    x_0^{-a_0}\cdots x_n^{-a_n},\qquad a_i\ge1,\quad
    \sum_i a_i=-m,
\]
on the full intersection of the standard affine cover.
\end{theorem}

\begin{proof}
Use the standard affine cover \(U_i=D_+(x_i)\).  Since all finite intersections
are affine and \(\OO(m)\) is quasi-coherent, Cech cohomology for this cover
computes sheaf cohomology.  A section over
\[
    U_{i_0}\cap\cdots\cap U_{i_p}
\]
is the degree \(m\) part of
\[
    A[x_0,\ldots,x_n]_{x_{i_0}\cdots x_{i_p}}.
\]
Thus the Cech complex is the degree \(m\) part of the localization complex
\[
0\to S\to\prod_i S_{x_i}\to
\prod_{i<j}S_{x_ix_j}\to\cdots\to S_{x_0\cdots x_n}\to0,
\]
where \(S=A[x_0,\ldots,x_n]\).  Decompose this complex as a direct sum over
Laurent monomials \(x^u=x_0^{u_0}\cdots x_n^{u_n}\) of total degree \(m\).  For
a fixed exponent vector \(u\), the monomial appears in the term indexed by a
set \(I\) exactly when \(u_j\ge0\) for all \(j\notin I\).  Hence its
contribution is the augmented cochain complex of a simplex on the set of
indices where \(u_j<0\).  This complex is exact unless all \(u_j<0\); if no
exponent is negative, the monomial contributes to \(H^0\).  Therefore \(H^0\)
consists of ordinary homogeneous polynomials of degree \(m\), and the only
possible higher cohomology occurs in top degree \(n\), represented by monomials
with all exponents negative.
\end{proof}

\begin{corollary}
If \(k\) is a field, then the cohomology groups
\[
    H^q(\PP^n_k,\OO(m))
\]
are finite-dimensional \(k\)-vector spaces for all \(q,m\), and vanish for
\(q>n\).
\end{corollary}

\begin{proof}
The theorem gives an explicit finite basis in degree \(0\) for \(m\ge0\), no
degree \(0\) sections for \(m<0\), vanishing in intermediate degrees, and in top
degree a finite basis of Laurent monomials because the equations \(a_i\ge1\)
and \(\sum_i a_i=-m\) have only finitely many solutions.  Cech complexes for
the standard cover have length \(n\), so cohomology vanishes for \(q>n\).
\end{proof}

\section{Serre Vanishing and Projective Finiteness}

\begin{theorem}[Serre vanishing]
Let \(A\) be a noetherian ring, let \(X=\PP^n_A\), and let \(\FF\) be a
coherent sheaf on \(X\).  Then there exists \(N\) such that for all \(m\ge N\),
\[
    H^q(X,\FF(m))=0\qquad(q>0).
\]
\end{theorem}

\begin{proof}
Every coherent sheaf on \(\PP^n_A\) is the cokernel of a morphism
\[
    \bigoplus_{\alpha}\OO(-r_\alpha)\longrightarrow
    \bigoplus_{\beta}\OO(-s_\beta)
    \longrightarrow \FF\longrightarrow0
\]
with finitely many summands.  To see this, cover \(X\) by standard affines,
choose finitely many local generators, multiply by sufficiently large powers of
the relevant \(x_i\) to extend them to global sections of twists, and use
coherence to make the kernel finitely generated.  Twisting by \(m\) gives
\[
    \bigoplus_{\alpha}\OO(m-r_\alpha)\to
    \bigoplus_{\beta}\OO(m-s_\beta)\to \FF(m)\to0.
\]
For \(m\) large, the higher cohomology of the two sums of line bundles vanishes
by the computation of projective space cohomology.  The long exact sequence
then gives \(H^q(X,\FF(m))=0\) for all \(q>0\), after increasing \(m\) and
inducting on \(q\) through the finite presentation just constructed.
\end{proof}

\begin{theorem}[Serre global generation]
Under the hypotheses of Serre vanishing, there exists \(N\) such that
\(\FF(m)\) is generated by global sections for all \(m\ge N\).
\end{theorem}

\begin{proof}
Choose a finite presentation
\[
    \bigoplus_{\alpha}\OO(-r_\alpha)\to
    \bigoplus_{\beta}\OO(-s_\beta)\to \FF\to0.
\]
For \(m\ge\max_\beta s_\beta\), each \(\OO(m-s_\beta)\) is generated by
homogeneous polynomials of degree \(m-s_\beta\), hence by global sections.  The
direct sum is globally generated, and its quotient \(\FF(m)\) is globally
generated.
\end{proof}

\begin{theorem}[Coherent finiteness for projective schemes]
Let \(k\) be a field, \(X\) a projective \(k\)-scheme, and \(\FF\) a coherent
sheaf on \(X\).  Then \(H^q(X,\FF)\) is finite-dimensional over \(k\) for all
q, and \(H^q(X,\FF)=0\) for \(q>\dim X\).
\end{theorem}

\begin{proof}
Choose a closed immersion \(i:X\hookrightarrow\PP^n_k\).  Since \(i\) is a
closed immersion, \(i_*\) is exact on coherent sheaves and
\[
    H^q(X,\FF)=H^q(\PP^n_k,i_*\FF).
\]
Thus it suffices to work on projective space.  For \(m\gg0\), Serre vanishing
gives no higher cohomology, and \(H^0(\FF(m))\) is a quotient of a finite direct
sum of \(H^0(\OO(r))\), hence finite-dimensional.  Choose a hyperplane
\(H\subset\PP^n_k\) not containing any associated component of \(\FF\).
Multiplication by its defining section gives an exact sequence
\[
    0\to\FF(m-1)\to\FF(m)\to \FF_H(m)\to0.
\]
The sheaf \(\FF_H\) is coherent on \(H\simeq\PP^{n-1}_k\).  By induction on
\(n\), the cohomology of \(\FF_H(m)\) is finite-dimensional.  The long exact
sequence then propagates finiteness from \(m\) to \(m-1\).  Repeating finitely
many times reaches \(m=0\).  The same hyperplane induction shows vanishing
above \(\dim X\), because after restricting to hyperplanes the dimension drops
and the long exact sequence has no non-zero terms beyond that bound.
\end{proof}

\begin{definition}
For a coherent sheaf \(\FF\) on a projective \(k\)-scheme \(X\) with a chosen
very ample sheaf \(\OO_X(1)\), the function
\[
    P_\FF(m)=\chi(X,\FF(m))
\]
is called the \emph{Hilbert function} of \(\FF\).
\end{definition}

\begin{theorem}[Hilbert polynomial]
For \(m\gg0\), the function \(P_\FF(m)\) agrees with a polynomial in \(m\) with
rational coefficients.  If \(X\) has dimension \(d\), this polynomial has
degree at most \(d\).
\end{theorem}

\begin{proof}
Embed \(X\) in projective space and replace \(\FF\) by its direct image.  By
Serre vanishing, for \(m\gg0\),
\[
    P_\FF(m)=h^0(X,\FF(m)).
\]
Choose a finite graded \(S=k[x_0,\ldots,x_n]\)-module \(M\) with
\(\widetilde M\simeq\FF\).  For \(m\gg0\), \(H^0(X,\FF(m))\simeq M_m\).  The
Hilbert-Serre theorem for finite graded modules over a polynomial ring states
that \(\dim_k M_m\) agrees for large \(m\) with a polynomial whose degree is the
dimension of the support of \(M\).  This is a pure commutative algebra theorem,
and gives the result.
\end{proof}

\section{Exercises Used Later}

\begin{exercise}[Cech \(H^1\) on a two-open cover]\label{ex:cech-two-open}
Let \(X=U\cup V\) and let \(\FF\) be a sheaf of abelian groups.  Write the Cech
complex for this cover and identify \(\cH^1(\{U,V\},\FF)\).
\end{exercise}

\begin{solution}
The complex is
\[
    \FF(U)\oplus\FF(V)\longrightarrow \FF(U\cap V),
    \qquad (s,t)\longmapsto s|_{U\cap V}-t|_{U\cap V}.
\]
There are no higher intersections.  Hence
\[
    \cH^1(\{U,V\},\FF)=
    \FF(U\cap V)/
    \{s|_{U\cap V}-t|_{U\cap V}:s\in\FF(U),\,t\in\FF(V)\}.
\]
\end{solution}

\begin{exercise}[Affine acyclicity for principal covers]\label{ex:principal-acyclic}
Let \(A\) be a ring, \(M\) an \(A\)-module, and \(f,g\in A\) with
\((f,g)=A\).  Prove exactness of
\[
    0\to M\to M_f\oplus M_g\to M_{fg}\to0.
\]
\end{exercise}

\begin{solution}
The first map sends \(m\) to \((m,m)\), and the second sends \((a,b)\) to
\(a-b\).  If \(m\) maps to zero, then \(m=0\) after localizing at \(f\) and at
\(g\); since \(af+bg=1\) for some \(a,b\), a sufficiently large power argument
gives \(m=0\).  If \((u,v)\) maps to zero in \(M_{fg}\), then \(u\) and \(v\)
agree after localizing at \(fg\).  After clearing denominators, choose \(N\)
such that \(f^Ng^N(u-v)=0\).  A partition-of-unity calculation using
\((af+bg)^N=1\) produces an element \(w\in M\) restricting to \(u\) in \(M_f\)
and to \(v\) in \(M_g\).  Finally, every element of \(M_{fg}\) is locally a
difference of elements from \(M_f\) and \(M_g\), so the last map is surjective.
\end{solution}

\begin{exercise}[Euler characteristic of a skyscraper]\label{ex:skyscraper-euler}
Let \(X\) be a noetherian scheme over \(k\), and let \(P\) be a closed point.
Show that the skyscraper sheaf \(\kappa(P)\) has
\[
    \chi(X,\kappa(P))=[\kappa(P):k].
\]
\end{exercise}

\begin{solution}
The sheaf is supported at one closed point.  Its global sections are
\(\kappa(P)\), because every open set containing \(P\) sees the same value and
opens not containing \(P\) see zero.  It is flasque: restrictions are either
the identity, the map to zero, or zero.  Flasque sheaves have no higher
cohomology in the Godement definition.  Thus the Euler characteristic is
\(\dim_k\kappa(P)\).
\end{solution}

\begin{exercise}[Global generation gives a map to projective space]\label{ex:global-generation-map}
Let \(\LL\) be an invertible sheaf on a scheme \(X\), generated by global
sections \(s_0,\ldots,s_n\).  Show that these sections define a morphism
\[
    X\longrightarrow \PP^n
\]
on the open cover where some \(s_i\) is non-zero.
\end{exercise}

\begin{solution}
On the open set \(U_i=\{s_i\ne0\}\), trivialize \(\LL\) by \(s_i\).  Then the
ratios \(s_j/s_i\) are regular functions on \(U_i\), giving a morphism
\[
    U_i\to D_+(x_i)\simeq\AA^n.
\]
On \(U_i\cap U_j\), the two sets of affine coordinates are related by the
standard transition functions of projective space.  Hence the local morphisms
glue to a morphism \(X\to\PP^n\).
\end{solution}

\section{Higher Direct Images}

\begin{definition}
Let \(f:X\to Y\) be a continuous map and let \(\FF\) be a sheaf of abelian
groups on \(X\).  The \(i\)-th \emph{higher direct image} \(R^if_*\FF\) is the
sheaf associated to the presheaf
\[
    V\longmapsto H^i(f^{-1}V,\FF|_{f^{-1}V})
\]
on open subsets \(V\subset Y\).
\end{definition}

\begin{proposition}
There is a natural identification
\[
    R^0f_*\FF=f_*\FF.
\]
If
\[
    0\to\FF'\to\FF\to\FF''\to0
\]
is exact, then there is a long exact sequence of higher direct images.
\end{proposition}

\begin{proof}
For \(i=0\), the defining presheaf is
\[
    V\mapsto H^0(f^{-1}V,\FF)=\Gamma(f^{-1}V,\FF)=f_*\FF(V),
\]
which is already a sheaf.  For an exact sequence, apply the long exact
cohomology sequence to each open set \(f^{-1}V\).  The resulting exact sequence
of presheaves becomes exact after sheafification because exactness of sheaves
can be checked on stalks.
\end{proof}

\begin{proposition}
If \(Y=\Spec A\) is affine, then
\[
    H^i(Y,R^0f_*\FF)=H^i(Y,f_*\FF)
\]
and the edge map
\[
    H^i(Y,f_*\FF)\to H^i(X,\FF)
\]
is an isomorphism when \(f\) is affine and \(\FF\) is quasi-coherent.
\end{proposition}

\begin{proof}
If \(f\) is affine and \(\FF\) is quasi-coherent, then \(f_*\FF\) is
quasi-coherent.  For every affine open \(V\subset Y\), the inverse image
\(f^{-1}V\) is affine, so \(H^q(f^{-1}V,\FF)=0\) for \(q>0\).  Thus
\[
    R^qf_*\FF=0\qquad(q>0).
\]
Computing cohomology by an affine Cech cover of \(Y\), the Cech complex for
\(f_*\FF\) is the same as the Cech complex for \(\FF\) on the inverse-image
cover of \(X\).  Hence the edge map is an isomorphism.
\end{proof}

\begin{theorem}[Coherence of higher direct images]
Let \(f:X\to Y\) be a projective morphism of noetherian schemes and let \(\FF\)
be coherent on \(X\).  Then \(R^if_*\FF\) is coherent on \(Y\) for every
\(i\ge0\), and \(R^if_*\FF=0\) for \(i\) sufficiently large.
\end{theorem}

\begin{proof}
The assertion is local on \(Y\), so assume \(Y=\Spec A\) and
\(X\hookrightarrow\PP^n_A\).  Replacing \(\FF\) by its direct image on
\(\PP^n_A\), use a finite presentation by sums of twists \(\OO(m)\).  The
cohomology of these twists over projective space is computed by the explicit
Cech complex of the standard affine cover; each term of that complex is a
finite \(A\)-module in the relevant degrees.  Therefore the cohomology modules
are finite \(A\)-modules.  These modules are the sections of \(R^if_*\FF\) over
\(\Spec A\), and localization of the Cech complex identifies their restrictions
to principal opens.  Hence the associated sheaves are coherent.  Vanishing in
large degree follows from the finite length of the Cech complex.
\end{proof}

\begin{theorem}[Cohomology and base change, curve form]
Let \(f:X\to S\) be a projective flat morphism of noetherian schemes whose
fibers are curves, and let \(\FF\) be coherent and flat over \(S\).  For
\(s\in S\), there is a natural map
\[
    (R^if_*\FF)\tensor_{\OO_S}\kappa(s)
    \longrightarrow
    H^i(X_s,\FF_s).
\]
It is an isomorphism for \(i=0\) when the dimension of
\(H^0(X_s,\FF_s)\) is locally constant near \(s\), and for \(i=1\) under the
corresponding local constancy of \(H^1\).
\end{theorem}

\begin{proof}
Choose a finite affine open cover of \(X\) over an affine neighbourhood of
\(s\).  Because the fibers have dimension one, the Cech complex computing the
cohomology of \(\FF\) has only two non-zero degrees after refinement.  Flatness
of \(\FF\) over \(S\) makes the terms of this Cech complex flat modules over
the base.  Tensoring the complex with \(\kappa(s)\) gives the Cech complex of
\(\FF_s\) on the fiber.  The base change map is the natural map from cohomology
of a complex tensored with \(\kappa(s)\) to the cohomology of the tensor
complex.  For a two-term complex of finite flat modules, this map is an
isomorphism exactly when the relevant cokernel or kernel has locally constant
rank; this is the elementary linear algebra form of cohomology and base
change.  Applying it in degrees \(0\) and \(1\) gives the result.
\end{proof}

\begin{exercise}[Affine morphisms have no higher direct images]\label{ex:affine-higher-direct}
Let \(f:X\to Y\) be affine and let \(\FF\) be quasi-coherent.  Show that
\[
    R^if_*\FF=0\qquad(i>0).
\]
\end{exercise}

\begin{solution}
For every affine open \(V\subset Y\), the inverse image \(f^{-1}V\) is affine.
The restriction of \(\FF\) to this inverse image is quasi-coherent, so affine
acyclicity gives
\[
    H^i(f^{-1}V,\FF)=0
\]
for \(i>0\).  The presheaf defining \(R^if_*\FF\) is therefore zero on an
affine basis of \(Y\), hence its sheafification is zero.
\end{solution}
